\section{Pesquisas}

\section{Setup do Ambiente de Desenvolvimento}
    \subsection{IDE Utilizada}
    \subsection{Github e Controle de Versão}
        % \subsubsection{Branchs}

\section{Arquitetura do Software} \label{sec:arquitetura}

O aplicativo foi desenvolvido com o \textit{framework} Flutter e segue a estrutura de diretórios padrão de projetos Dart voltados a dispositivos móveis~\cite{flutter_docs}. Essa organização separa a configuração nativa da plataforma, os recursos estáticos embarcados, as ferramentas auxiliares de linha de comando, os testes automatizados e o código-fonte da aplicação. Na raiz do repositório, os principais arquivos e diretórios cumprem os seguintes papéis:

\begin{itemize}
    \item \texttt{pubspec.yaml} e \texttt{analysis\_options.yaml}: arquivos de configuração do projeto. O \texttt{pubspec.yaml} declara as dependências externas --- \texttt{camera}, \texttt{google\_\allowbreak mlkit\_\allowbreak face\_\allowbreak mesh\_\allowbreak detection}, \texttt{tflite\_\allowbreak flutter}, \texttt{sherpa\_\allowbreak onnx}, \texttt{sherpa\_\allowbreak onnx\_\allowbreak android\_\allowbreak arm64}, \texttt{audioplayers} e \texttt{path\_\allowbreak provider} --- e lista os diretórios de recursos estáticos empacotados na compilação, enquanto o \texttt{analysis\_\allowbreak options.yaml} ativa as regras de verificação estática do pacote \texttt{flutter\_lints}.
    \item \texttt{android/}: projeto nativo Android gerenciado pelo Gradle, que hospeda o motor do Flutter em uma \texttt{Activity} (Seção~\ref{sec:android}), define a compilação para a arquitetura \texttt{arm64-v8a} e declara no manifesto (\texttt{AndroidManifest.xml}) a permissão de uso da câmera frontal (\texttt{android.\allowbreak permission.\allowbreak CAMERA}).
    \item \texttt{assets/}: recursos estáticos incluídos no pacote de instalação (\texttt{APK}) para permitir o funcionamento sem conexão com a internet. Divide-se em \texttt{assets/phrases/}, com os cinco arquivos JSON do \textit{corpus} de frases de \gls{CAA} (Seção~\ref{sec:met_corpus}), e \texttt{assets/models/}, que armazena o modelo de detecção de pupila (\texttt{pupil\_128.tflite}), os parâmetros geométricos do recorte ocular (\texttt{crop\_params.json}) e o diretório \texttt{kokoro\_pt/} com o modelo de síntese de voz (\texttt{model.onnx}), a tabela de vozes (\texttt{voices.bin}), o vocabulário fonético (\texttt{tokens.txt}), o manifesto de integridade (\texttt{manifest.json}) e os dados fonéticos do \texttt{espeak-ng} para o português brasileiro.
    \item \texttt{tool/}: ferramentas executadas na máquina de desenvolvimento, fora do aplicativo. Contém o \textit{script} \texttt{fetch\_kokoro\_model.py}, que baixa o pacote de voz (cerca de 355~MB, mantido fora do controle de versão), extrai os arquivos necessários ao português brasileiro e confere o tamanho e o resumo SHA-256 de cada arquivo; o programa \texttt{evaluate\_predictor.dart}, que executa a avaliação offline do preenchimento automático (Seção~\ref{sec:met_avaliacao}); e o \textit{script} \texttt{render\_facemesh.py}, usado para gerar figuras da malha facial.
    \item \texttt{test/}: suíte de testes automatizados de unidade e de interface, executada com \texttt{flutter test}. Os testes cobrem a geometria do recorte ocular, a leitura de quadros \texttt{NV21}, o cálculo de rotação da câmera, o ajuste polinomial da calibração, a predição de palavras, o cálculo das métricas e o controle de estado do sintetizador de voz, permitindo verificar a lógica dos módulos sem depender do \textit{tablet} a cada alteração.
    \item \texttt{lib/}: diretório que reúne todo o código-fonte em Dart do aplicativo.
\end{itemize}

\section{Utilização de IA no Desenvolvimento}
